\documentclass[a4paper,12pt,oneside]{report}
\usepackage{graphicx}
\usepackage{color}
\usepackage{url}
\usepackage{subfigure}
\usepackage[utf8]{inputenc}
\usepackage[T1]{fontenc}
%\usepackage{tgpagella}
\usepackage{xstring}
\usepackage{wrapfig}
\usepackage{appendix}
\usepackage{nomencl}

\newcommand{\myscale}{0.74}
\newcommand{\vect}[1]{\boldsymbol{#1}}
\newcommand{\code}[1]{\texttt{\StrSubstitute{#1}{.}{.\.}}}
\def\.{\discretionary{}{}{}}
\newcommand{\jmodule}[1]{\texttt{\textit{#1}}}
\linespread{1.5}

\setlength{\hoffset}{-1in} %left margin will be 0, as hoffset is by default 1inch
\setlength{\voffset}{-1in} %analogous voffset
\setlength{\oddsidemargin}{4cm}
\setlength{\evensidemargin}{4cm}
\setlength{\topmargin}{25mm}
\setlength{\footskip}{1cm}
\setlength{\headheight}{0cm}
\setlength{\headsep}{0cm}
\setlength{\marginparwidth}{0cm}
\setlength{\marginparpush}{0cm}
\setlength{\textheight}{23.7cm}
\setlength{\textwidth}{14.5cm}
\let\openright=\clearpage

\def\mfauthor{Matej Vitásek}
\def\mfadvisor{RNDr. Irena Mlýnková Ph.D.}
\def\mfplacedate{Praha, 2011}

% Tato makra přesvědčují mírně ošklivým trikem LaTeX, aby hlavičky kapitol
% sázel příčetněji a nevynechával nad nimi spoustu místa. Směle ignorujte.
\makeatletter
\def\@makechapterhead#1{
  {\parindent \z@ \raggedright \normalfont
   \Huge\bfseries \thechapter. #1
   \par\nobreak
   \vskip 20\p@
}}
\def\@makeschapterhead#1{
  {\parindent \z@ \raggedright \normalfont
   \Huge\bfseries #1
   \par\nobreak
   \vskip 20\p@
}}
\makeatother

\def\chapwithtoc#1{
\chapter*{#1}
\addcontentsline{toc}{chapter}{#1}
}

\renewcommand{\nomname}{List of Abbreviations}
\makenomenclature

\begin{document}

\lefthyphenmin=2
\righthyphenmin=2

%%% Titulní strana práce ======================================================
\pagestyle{empty}
\begin{center}

\large

Charles University in Prague

\medskip

Faculty of Mathematics and Physics

\vfill

{\bf\Large MASTER THESIS}

\vfill

\centerline{\mbox{\includegraphics[width=60mm]{logo.eps}}}

\vfill
\vspace{5mm}

{\LARGE \mfauthor}

\vspace{15mm}

% exactly as assigned
{\LARGE\bfseries Inference of XML Integrity Constraints}

\vfill

% Název katedry nebo ústavu, kde byla práce oficiálně zadána
% (dle Organizační struktury MFF UK)
Department of Software Engineering

\vfill

\begin{tabular}{rl}

Supervisor of the master thesis: & 	\mfadvisor \\
\noalign{\vspace{2mm}}
Study programme: & Informatika \\
\noalign{\vspace{2mm}}
Specialization: & ISS \\
\end{tabular}

\vfill

% Zde doplňte rok
\mfplacedate 

\end{center}



\newpage % ============================================================
%%% Následuje vevázaný list -- kopie podepsaného "Zadání diplomové práce".
%%% Toto zadání NENÍ součástí elektronické verze práce, nescanovat.
\openright

\noindent
[Sample: Here you may thank whoever you wish (the supervisor of the thesis, the consultant,
the person who lent the software, literature, etc.)]

TODO thank jInfer team, Mlynkova, anyone helping create this, colleagues at work...


\newpage % ============================================================
%%% Strana s čestným prohlášením k diplomové práci
\vglue 0pt plus 1fill
\noindent
I declare that I carried out this master thesis independently, and only with the cited
sources, literature and other professional sources.

\medskip\noindent
I understand that my work relates to the rights and obligations under the Act No.
121/2000 Coll., the Copyright Act, as amended, in particular the fact that the Charles
University in Prague has the right to conclude a license agreement on the use of this
work as a school work pursuant to Section 60 paragraph 1 of the Copyright Act.

\vspace{10mm}

\hbox{\hbox to 0.5\hsize{%
In ........ date ............
\hss}\hbox to 0.5\hsize{%
signature
\hss}}

\vspace{20mm}



\newpage % ============================================================
%%% Povinná informační strana diplomové práce
\vbox to 0.5\vsize{
\setlength\parindent{0mm}
\setlength\parskip{5mm}

Název práce: Odvozování integritních omezení v XML % TODO how to translate this correctly?

Autor: \mfauthor

Katedra:  Katedra softwarového inženýrství

Vedoucí diplomové práce: \mfadvisor{}, Katedra softwarového inženýrství

Abstrakt: [abstract of 80-200 words in Czech, but not a copy of the assignment of the
master thesis]

Klíčová slova:  XML, ID atributy, odvozování

\vss}\nobreak\vbox to 0.49\vsize{
\setlength\parindent{0mm}
\setlength\parskip{5mm}

Title:  Inference of XML Integrity Constraints

Author: \mfauthor

Department: Department of Software Engineering

Supervisor: \mfadvisor{}, Department of Software Engineering

Abstract:  [abstract of 80-200 words in English, but not a copy of the assignment of
the master thesis]

Keywords: XML, ID attributes, inference

\vss}

\newpage

%%% Strana s automaticky generovaným obsahem diplomové práce. U matematických
%%% prací je přípustné, aby seznam tabulek a zkratek, existují-li, byl umístěn
%%% na začátku práce, místo na jejím konci.

\openright
\pagestyle{plain}
\setcounter{page}{1}
\tableofcontents

% =============== TEXT ================
\newpage

% TODO what style of quotation marks to use? 
% TODO how to format text? is textit / texttt / ... OK?

\chapwithtoc{Preface}

Along with technologies like JSON, SQL/noSQL databases and bla, XML is one of the leading formats for storing structured data. However, even though languages such as DTD and XML Schema to describe XML structure already exist since a long time, most of the documents use outdated or no schema at all (link Vosta's Even ant can create...). To tackle this problem one may employ reverse-engineering techniques to infer the schema from existing documents, such as those described in A, B, C, jInfer. But the schema is not the only constraint that can be imposed on an XML document: the concept of \textit{keys} and \textit{foreign keys}, well known from the relational database world, applies here as well. One could go even further and try to find even more sophisticated relations in the data, such as \textit{functional dependencies} (link Sviro).

This work will be building upon the achievements of jInfer schema inference framework, expanding its possibilities in the field of search for \textit{key-} and \textit{foreign key-}like structures in existing XML documents.

TODO Integrity constraints can be keys (with ID attributes as a "sub-group"), FKs, functional dependencies (quote Sviro), etc.
We will focus on the first kind - ID and IDREF attributes.

TODO argument: test data with DTD (and thus possibly ID/IDREF) is more common, we can have better test sets.

\section{Structure of the thesis}

The thesis will be structured as follows. 

First, we will introduce a few notions required throughout the work, such as XML tree, ID attributes, ID sets and keys for XML. 

Secondly, we will review approaches to ID attribute and XML keys search from previous articles on this topic. 

This will lead us to the NP-complete problem of maximal independent set (IS)\nomenclature{IS}{Independent Set}, where we will inspect approaches for solving it.

We will discuss a closely related Mixed Integer Problem (MIP)\nomenclature{MIP}{Mixed Integer Problem} and prove their "equality".

Afterwards, we will show how to use an external MIP solver and various heuristics to tackle this problem.

An extension to jInfer for finding ID attributes using MIP solver and a combination of heuristics will be presented and experimentally evaluated in the closing chapters.

\chapter{Definitions}

\begin{itemize}
	\item XML tree - from FIDAX or wherever
	\item Element, Attribute
	\item ID, IDREF, IDREFS attribute according to specification
	\item XSD keys (compare them to ID attributes)
	\item Key according to Keys for XML \cite{keX}
	\item Attribute mapping
	\item ID set
  \item Weight: support, coverage
\end{itemize}

\chapter{Research}

According to the article \cite{fidax}, ...

TODO talk about FIDAX
To the best of our knowledge, there are no other articles dealing with this problem. 

TODO talk about Fajt - but that's different, that's keys
TODO add citation for Fajt

\section{Independent Set}

TODO define IS + MIS rigorously

TODO other approaches to max IS

\chapter{MIP Approach}

TODO rigorously define MIP - optimization problem, some of the variables integral

TODO FIDAX is cool, but their heuristic might not be the best. It is a IS problem, let's try a multi-purpose solver. Use GLPK. \nomenclature{GLPK}{GNU Linear Programming Kit}

\section{Finding ID sets with GLPK}

TODO how GLPK works (branch&bound), how to construct a problem, how output looks, how to interpret it, limit time-> limit quality, ...
 
 -> wow, GLPK works. However, it takes too long to get to optimum (link experiments), so let's try other heuristics.
 
\section{Heuristics}

TODO what is heuristics (link wise books), what is metaheuristics (we will be using them)

TODO mention things like Taboo search and Genetic Algorithms (we can emulate them with Crossover/Mutation) (we won't be using them)

TODO we will be working with heuristics striving to do the following: input is a list of AMs, they have their weight, we try to find a non-conflicting subset which maximizes the weight

\subsection{Constructions Heuristics}

% TODO can I format things in nomenclature?
\nomenclature{CH}{Construction Heuristic} %TODO decide on how to format CH/IH 

TODO construction heuristics are heus that provide us with at least some solution.

\subsubsection{FIDAX}

TODO this is a heuristic too! We can start with it!

\subsubsection{Random}

TODO this is trivial, but works - link experiments where it bashes FIDAX

\subsubsection{Fuzzy}

TODO kind of trivial

\subsubsection{Incremental}

TODO trivial, hungry

\subsubsection{Removal}

TODO trivial, "hungry"

\subsubsection{Truncated Branch \& Bound}

TODO if we limit GLPK's runtime, we get this 

\subsection{Improvement Heuristics}

\nomenclature{IH}{Improvement Heuristic}

TODO what they are, that they need a pool sometimes, how do we give them pool, ...

TODO mention intensification, diversification

\subsubsection{Identity}

TODO trivial trivial

\subsubsection{Remove Worst}

TODO trivial

\subsubsection{Random Remove}

TODO trivial, perturbation

\subsubsection{Hungry}

TODO trivial

\subsubsection{Mutation}

TODO

\subsubsection{Crossover}

TODO

\subsubsection{Local Branching}

TODO explain how this translates to GLPK input

\chapter{Experiments}

TODO describe all the inputs (take it from the mindmap), describe how the output looks (sample trace - link to appendix)

\section{Experimental Data}

TODO types of data:
\begin{itemize}
	\item Realistic 
  \item Realistic with artificial attributes
	\item Artificial
\end{itemize}

TODO for all types, show GraphViz pictures.

\subsection{Realistic data}

TODO talk about all that are used - mention the dataset name, some characteristics (excel sheet)

\subsection{Realistic data with artificial attributes}

TODO we have taken some of the text data and converted it to attributes like this, A B C

\subsection{Artificial data}

TODO these are random graphs.
TODO talk about how they are created, about their structure.
TODO talk about some "density".

\section{Experimental Setup}

TODO we will be using jInfer and its module AttributeStats see \ref{jInfer} and \ref{attrstats}.

\section{Experimental Results}

TODO following experiments + more

\begin{itemize}
  \item Which construction heuristic is the best - limit GLPK to like 1 second
  \item With the best construction heuristic, which IH is the best? Note that this needn't be the best combination!
  \item With the best IH from the previous step, how about using Random as CH?
	\item GLPK: Input -> How long the search for optimum lasts
	\item GLPK: Time limit -> Quality
	\item Structure of the data -> Which algorithms perform the best
	\item Alpha, Beta -> Quality
  \item Alpha, Beta -> How does the solution work?
	\item Which GLPK -> How does (time->quality) look, how does (input->time) look
	\item Do I beat FIDAX with Random?
	\item Does Fuzzy beat Random? (Seems that Random is better and even has smaller variance)
\end{itemize}

\section{The "Best" Algorithm}

TODO we have asked a lot of questions and got the answers, now is the time to summarize, to find some "wisdom".

First of all: if we have the time, it is best to just let the GLPK run. 
We will find the optimum in this way. 
And for many purposes, this is just fine - we need to infer something about the schema, we do it only once, so it doesn't matter that much how long it takes.

Second: if we don't have enough time, we should do XYZ but it might depend on the characteristics of the data.
For realistic (looks like THIS in graph representation) it is best to do A B C.
Whereas, for artificial data (having IJK as their representation) it is best to just do 1 2 3.

\chapter{Future work}

TODO extension to more documents, IDREFs, ...  

\chapwithtoc{Conclusion}

TODO conclusion: From all integrity constraints in XML we chose the ID attributes and decided to improve the search for them. We discussed the approach chosen in FIDAX. Based on this article we introduced the MIP approach and demonstrated how to find the optimal (w.r.t. weight) ID set.

However, this approach takes too long for some inputs, so we introduced a whole spectrum of construction as well as improvement heuristics. By combining these algorithms and tuning their parameters we were able to find very good solutions while maintaining low running times.

%%% Seznam použité literatury
\newpage
\nocite{*}
\bibliographystyle{alpha}
\bibliography{literature}

%%% Tabulky v diplomové práci, existují-li.
%\chapwithtoc{List of Tables}
\listoftables
\addcontentsline{toc}{chapter}{List of Tables}

%%% Použité zkratky v diplomové práci, existují-li, včetně jejich vysvětlení.
%\chapwithtoc{List of Abbreviations}

\printnomenclature[2cm]
\addcontentsline{toc}{chapter}{List of Abbreviations}

%%% Přílohy k diplomové práci, existují-li (různé dodatky jako výpisy programů,
%%% diagramy apod.). Každá příloha musí být alespoň jednou odkazována z vlastního
%%% textu práce. Přílohy se číslují.

\appendix

\openright
\addcontentsline{toc}{chapter}{Appendices}

\chapter{jInfer}
\label{jInfer}

TODO what it is, who created it and why, what it can do, where to get it.

\chapter{AttributeStats}
\label{attrstats}

TODO talk a bit about the structure of AttributeStats module.

\chapter{Traces}
\label{traces}

TODO a few (commented!) traces from experiment runs

\end{document}
